\subsection{Блоки и нитки}
%%%%%%%%%%%%%%%%%%%%%%%%%%%%%%%%%%%%%%%%%%%%%%%%%%%%
%%%%%%%%%%%%%%%%%%%%%%%%%%%%%%%%%%%%%%%%%%%%%%%%%%%%
\z{В какой из приведённых систем возможно равновесие грузов и при каком соотношении масс $M$ и $m$?
		\begin{center}%
 			\includegraphics[width=0.9\textwidth]{3.3.1}%
		\end{center}%
}
%
{\puan $M=2m$;~\puan НЕТ;~\puan $M=\frac{m}{2}$;~\puan НЕТ}

%%%%%%%%%%%%%%%%%%%%%%%%%%%%%%%%%%%%%%%%%%%%%%%%%%%%
\z{Груз массой $m = 3\, \text{кг}$ подвешен на нити к потолку лифта, который движется с ускорением $2g/3$. Найдите силу натяжения нити $T$, если лифт ускоряется:
\pu вверх;
\pu вниз.
}
%
{\puan $T=\frac{5}{3}mg$;~\puan $T=\frac{1}{3}mg$}
%%%%%%%%%%%%%%%%%%%%%%%%%%%%%%%%%%%%%%%%%%%%%%%%%%%%
\z[{\begin{wrapfigure}[4]{r}{0.15\textwidth}%
		\vspace{-45pt}%
		\begin{flushright}%
 			\includegraphics[width=0.15\textwidth]{3.3.3}%
		\end{flushright}%
	\end{wrapfigure}%
	}]%
{Мальчик соорудил себе подвесное кресло и с его помощью смог приподнять себя над землёй. С какой силой для этого мальчик должен был тянуть верёвку, и с какой силой он давил на сиденье, если его масса $M = 50 \ \text{кг}$, а масса кресла $m = 10 \ \text{кг}$?}
%
{$T = \frac{2(M + m)g}{3} = 400 \, \text{Н}, \quad N = \frac{M - 2m}{3} g = 100 \, \text{Н}$}

%%%%%%%%%%%%%%%%%%%%%%%%%%%%%%%%%%%%%%%%%%%%%%%%%%%%
\z[{\begin{wrapfigure}[3]{r}{0.3\textwidth}%
		\vspace{-10pt}%
		\begin{flushright}%
 			\includegraphics[width=0.3\textwidth]{3.3.4}%
		\end{flushright}%
	\end{wrapfigure}%
	}]%
{Три груза массой $4m$, $3m$ и $m$, связанные невесомыми нерастяжимыми нитями, движутся по гладкой горизонтальной поверхности под действием горизонтальной силы $F$, приложенной к грузу массой $m$. Найдите ускорение, с которым движутся грузы и силы натяжения обеих нитей.}
%
{$a=\frac{F}{8m};~T_1=\frac{F}{2};~T_2=\frac{7F}{8}$}


%%%%%%%%%%%%%%%%%%%%%%%%%%%%%%%%%%%%%%%%%%%%%%%%%%%%
\z[{\begin{wrapfigure}[3]{r}{0.3\textwidth}%
		\vspace{-10pt}%
		\begin{flushright}%
 			\includegraphics[width=0.3\textwidth]{3.3.5}%
		\end{flushright}%
	\end{wrapfigure}%
	}]%
{При каком минимальном значении горизонтальной силы $F$ разорвется связывающая грузы легкая нить, если она выдерживает предельную нагрузку $T_0$?  Массы грузов $2m$ и $5m$. Трением пренебречь.}
%
{при $F>\frac{7}{2}T_0$}


%%%%%%%%%%%%%%%%%%%%%%%%%%%%%%%%%%%%%%%%%%%%%%%%%%%%
\z[{\begin{wrapfigure}[3]{r}{0.3\textwidth}%
		\vspace{-10pt}%
		\begin{flushright}%
 			\includegraphics[width=0.3\textwidth]{3.3.6}%
		\end{flushright}%
	\end{wrapfigure}%
	}]%
{На гладком столе лежат два связанных нитью груза массой $m_1=200\, \text{г}$ и $m_2=300 \, \text{г}$. К левому грузу приложена $F_1=0.6 \, \text{Н}$, к правому -- $F_2=0.1\, \text{Н}$. С каким ускорением будут двигаться грузы? Чему равна сила натяжения нити? }
%
{$a=\frac{F_1-F_2}{m_1+m_2}=1\,\text{м/с}^2;~T=\frac{m_1F_2+m_2F_1}{m_1+m_2}=0.4\,\text{Н}$}




%%%%%%%%%%%%%%%%%%%%%%%%%%%%%%%%%%%%%%%%%%%%%%%%%%%%
\z[{\begin{wrapfigure}[3]{r}{0.3\textwidth}%
		\vspace{-10pt}%
		\begin{flushright}%
 			\includegraphics[width=0.3\textwidth]{3.3.7}%
		\end{flushright}%
	\end{wrapfigure}%
	}]%
{Какую горизонтальную силу надо прикладывать к тележке, чтобы удержать ее на месте? Масса груза $m$ известна.}
%
{$F=3mg$}